\section{Introduction}

The useful model for a sentiment project depends on how many labels exist and what deployment costs matter. A linear classifier over sparse TF-IDF features can be accurate with little training data, while a neural sequence model must learn both token representations and a decision boundary. A single full-data leaderboard misses the point at which that trade-off changes.

Part~1 is the approved project contribution. It benchmarks nine complete pipelines across three Korean review domains, beginning at 200 labels and ending at each cleaned training pool. The primary question is whether the best learned-from-scratch neural method exceeds the best TF-IDF method by at least 0.01 mean macro-F1, and whether that crossover persists at every larger tested size. We pair this performance question with measured fit time, inference latency, memory, and serialized size.

Part~2 is a separate bonus extension. It asks whether 4,000 labels already available in two source domains can reduce the number of new labels required in the third. The comparison first isolates the effect of pooling, then tests standard supervised contrastive learning and a domain-balanced cross-domain variant. Part~1 TF-IDF and target-only Transformer results are reused rather than retrained.

The study contributes:
\begin{itemize}
  \item a controlled three-domain scale benchmark of four classical and five learned-from-scratch neural pipelines;
  \item a joint analysis of macro-F1 crossover and computational cost; and
  \item a bonus evaluation of domain-balanced cross-domain supervised contrastive learning under limited target labels.
\end{itemize}
The experiments do not cover continual learning or pretrained Korean language models.
